\section{Conclusion and Outlook}
\label{sec:conclusion}

We have presented the first nucleon gluon PDF calculated through LaMET with hybrid-ratio renormalization with at two heavier than physical pion masses, $690$ and $310$ MeV, using two gauge link smearing techniques, Wilson flow with a flow time of $3a^2$ and 5 steps of hypercubic smearing.
Based on the prior work toward a LaMET gluon PDF~\cite{Good:2024iur}, we focused on a single gluon operator $O^{(3)}$ defined in Eq.~6 of Ref.~\cite{Good:2024iur}, which has been used in unpolarized gluon studies using pseudo-PDF method~\cite{Fan:2020cpa,Fan:2021bcr,HadStruc:2021wmh,Salas-Chavira:2021wui,Fan:2022kcb,Delmar:2023agv,Good:2023ecp}.
We computed the one-loop Wilson coefficient needed for hybrid renormalization for this operator and the corresponding one-loop matching kernel.
Using these Wilson coefficients, we found $\delta m+m_0$ to be consistent between the pion and nucleon, whether at the light or strange mass, in contrast with previous results for the $O^{(1)}$ and $O^{(2)}$ operators in Ref.~\cite{Good:2024iur}.
We computed the nucleon gluon PDFs for each pion mass and smearing technique with hadron boost momentum $P_z = 1.71$~GeV.
We found that the Wilson3 smearing significantly affects the physics, producing PDFs that are negative around $x \gtrsim 0.6$, while the noisier HYP5 PDFs do not fall significantly below zero.
We compared the HYP5 nucleon gluon PDFs at both pion masses with from the CT18~\cite{Hou:2019efy}, JAM24~\cite{Anderson:2024evk}, and CJ22~\cite{Cerutti:2025yji} global fits as a small sample.
However, we find good agreement with the CT18 PDF across the full region, and less agreement with CT22 and JAM24 around $x\approx 0.5$, which shows that our PDF falls comfortably within the range of those found in global fits.
\UPDATE{The final error bars on the lattice PDFs do not estimate systematics from the extrapolations to complete the Fourier transform, the demonstration of which was the central challenge in these first proof of principle results.
Bayesian and machine learning techniques~\cite{Chowdhury:2024ymm,Dutrieux:2024rem,Dutrieux:2025jed} could be transferred to expand on the simple extrapolation model to give a more rigorous error quantification.}
This study demonstrates a strong proof-of-principle for using large momentum effective theory to calculate gluon PDFs, which was long thought to be out of reach.
We believe that this work provides strong motivation to use time and resources to replace the large pion mass and smearing with noise-reduction techniques with fewer physics effects to obtain more reliable gluon PDFs from LaMET.
We plan to improve the current calculation by using self-renormalization scheme~\cite{LatticePartonLPC:2021gpi} with zero-momentum boosted matrix elements with multiple lattice spacings.
\UPDATE{This will allow us to better understand discretization and smearing effects and improve the systematics on the gluon PDFs by taking the continuum-physical limit.}
As the errors become smaller, there will also be a need to understand the effects of scale variation in the gluon PDFs from LaMET.
